\documentclass[../Main.tex]{subfiles}


% Definitions and setup

\onehalfspacing
\addbibresource{Source/sources.bib}


\begin{document}

\section{Einleitung}

\subsection{Motivation und Relevanz}

Hack-and-Slash Spiele gehören zu den technisch anspruchsvollsten Genres der Spieleentwicklung. Das Zusammenspiel von Spielerbewegung, Kampfsystem, Gegnerverhalten und Leveldesign muss genau aufeinander abgestimmt sein. Dabei ist das Fenster, in dem sich ein befriedigendes Spielerlebnis entfaltet, sehr eng. Genau dieses Zusammenspiel aber fasziniert uns als Spieler und stellt als Entwickler eine ebenso komplexe Herausforderung dar.

Die vorliegende Maturitätsarbeit entstand aus dem Wunsch, nicht nur Spiele zu spielen, sondern sie genauer zu verstehen und selbst eines zu entwickeln. Das Genre setzt dabei eine hohe Messlatte. Schon kleine Änderungen an einzelnen Systemen haben spürbare Auswirkungen auf das gesamte Spielgefühl. Dies macht die Entwicklung anspruchsvoll und bietet eine lehrreiche Erfahrung.

Da das Genre seine grösste Popularität zwischen 2001 und 2012 erlebte und heute nur noch eine kleine, aber treue Nische bedient, bietet es eine günstige Ausgangslage. Es existiert ausreichend Referenzmaterial aus einer designtechnisch reifen Ära, und gleichzeitig besteht Freiraum für moderne Interpretationen klassischer Mechaniken.


\subsection{Zielsetzung und Produktdefinition}

Ziel dieser Arbeit ist die Entwicklung eines spielbaren 3D-Hack-and-Slash Prototypen mit der Game Engine Godot. Der Schwerpunkt liegt dabei auf der technischen Umsetzung der zentralen Spielsysteme sowie deren Dokumentation und Evaluation.

Das Endprodukt umfasst einen zusammenhängenden Kampfabschnitt mit mindestens zwei unterscheidbaren Gegnertypen, die sich eigenständig bewegen und angreifen können. Die zentralen Spielmechaniken umfassen grundlegende Spielerbewegungen im dreidimensionalen Raum, Kampfmechaniken, Gegner-KI und übergeordnete Spielsysteme, Menüs oder Neustart nach dem Tod. Die Softwarearchitektur der wichtigsten Systeme wird durch Grafiken dokumentiert.

Ergänzend zur technischen Entwicklung werden zwei Referenzspiele des Genres anhand eines ausgearbeiteten Analyseframeworks untersucht. Die daraus gewonnenen Erkenntnisse fliessen direkt in das Design des Prototypen ein. Schliesslich wird der Prototyp in einem strukturierten Playtest mit ca. fünf Testpersonen evaluiert und die Ergebnisse ausgewertet.


\subsection{Erkenntnisfrage}

Im Zentrum unserer Arbeit steht folgende Erkenntnisfrage:

\begin{quote}
\textit{"Wie müssen Gegnerverhalten und Gegnerplatzierung in einem 3D-Hack-and-Slash-Prototyp gestaltet sein, damit ein zusammenhängender Kampfabschnitt mit nachvollziehbarem Spielfluss, angemessener Schwierigkeit und befriedigendem Spielergefühl entsteht?"}
\end{quote}

Diese Frage wird durch die Analyse bestehender Referenzspiele, die technische Umsetzung im Prototyp sowie die Auswertung der Playtests beantwortet.

\end{document}
